\documentclass[aspectratio=169]{beamer}
\usetheme{Madrid}
\usecolortheme{beaver}

\usepackage[UTF8,fontset=fandol]{ctex}
\usepackage{amsmath, amssymb}
\usepackage{booktabs}
\usepackage{ragged2e}
\usepackage{array}
\usepackage{tikz}
\usetikzlibrary{positioning, arrows.meta}

\newcolumntype{L}[1]{>{\RaggedRight\arraybackslash}p{#1}}

\title[0507report]{0507report}
\subtitle{}
\date{2026-05-07}

\begin{document}

\begin{frame}
    \titlepage
\end{frame}

% ============================================================
\section{当前结论}
% ============================================================

\begin{frame}[t]{主要修改内容}
\small
\begin{columns}[T]
\begin{column}{0.49\textwidth}
\textbf{Scope 扩散}
\begin{itemize}
    \item LLM 只判断 \texttt{scope\_amplitude}
    \item 本地代码用 embedding 算出的 task-family 相似度矩阵
    \item Gaussian 权重决定扩散到哪些相似任务
    \item 扩散扣的是其他任务的 \texttt{task\_self\_efficacy}
\end{itemize}
\end{column}

\begin{column}{0.49\textwidth}
\textbf{Stream memory}
\begin{itemize}
    \item 复用原生 \texttt{status + stream + search}
    \item 写入 raw episode 与检索相关经历
    \item 只服务 appraisal / attribution / interview
    \item 不新建第二套 memory 系统
\end{itemize}
\end{column}
\end{columns}

\vspace{0.2cm}
\end{frame}

% ============================================================
\section{Scope 扩散}
% ============================================================

\begin{frame}[t]{Scope 扩散在系统里的位置}
\small
\centering
\vspace{0.65cm}
\begin{tikzpicture}[
    node distance=0.72cm,
    box/.style={rectangle, rounded corners, draw=black!55, fill=black!4,
        align=center, minimum width=2.6cm, minimum height=0.72cm},
    arrow/.style={-{Latex[length=2mm]}, thick, draw=black!70}
]
\node[box] (event) {负面数字事件};
\node[box, right=of event] (attrib) {event attribution};
\node[box, right=of attrib] (amp) {\texttt{scope\_amplitude}};
\node[box, below=of amp] (gauss) {Gaussian spillover};
\node[box, left=of gauss] (se) {相似任务\\self-efficacy 下降};
\node[box, left=of se] (future) {后续任务\\更易回避/失败};

\draw[arrow] (event) -- (attrib);
\draw[arrow] (attrib) -- (amp);
\draw[arrow] (amp) -- (gauss);
\draw[arrow] (gauss) -- (se);
\draw[arrow] (se) -- (future);
\end{tikzpicture}

\vspace{0.45cm}
\begin{minipage}{0.78\textwidth}
\justifying
LLM 负责主观泛化强度，代码负责把这个强度按任务相似度扩散。
\end{minipage}
\end{frame}

\begin{frame}[t]{Scope 扩散公式与触发条件}
\small
\begin{columns}[T]
\begin{column}{0.52\textwidth}
\textbf{本地 Gaussian 公式}
\[
d(i,j)=1-S_{ij}
\]
\[
w_{ij}=\exp\left(-\frac{d(i,j)^2}{2\sigma^2}\right),\quad
\hat{w}_{ij}=\frac{w_{ij}}{\sum_{k\neq i}w_{ik}}
\]
\[
A_t=\beta\cdot U_t\cdot a_t,\quad
\Delta SE_{j,t}^{spill}=-A_t\hat{w}_{ij}
\]
\end{column}

\begin{column}{0.44\textwidth}
\textbf{默认参数}
\begin{itemize}
    \item \texttt{BETA=1.2}
    \item \texttt{THRESHOLD=0.15}
    \item \texttt{SIGMA=0.45}
\end{itemize}

\vspace{0.1cm}
\textbf{触发条件}
\begin{itemize}
    \item 负面 outcome
    \item attribution 状态为 \texttt{ok}
    \item \texttt{scope\_amplitude} 达到阈值
\end{itemize}
\end{column}
\end{columns}

\end{frame}

\begin{frame}[t]{embedding计算任务相似度矩阵}
\scriptsize
\centering
\renewcommand{\arraystretch}{1.06}
\begin{tabular}{@{}lcccccc@{}}
\toprule
 & Nav & Login & Search & Form & Service & Payment \\
\midrule
Nav & 1.000 & 0.635 & 0.692 & 0.620 & 0.746 & 0.603 \\
Login & 0.635 & 1.000 & 0.517 & 0.717 & 0.626 & 0.598 \\
Search & 0.692 & 0.517 & 1.000 & 0.621 & 0.711 & 0.693 \\
Form & 0.620 & 0.717 & 0.621 & 1.000 & 0.695 & 0.615 \\
Service & 0.746 & 0.626 & 0.711 & 0.695 & 1.000 & 0.643 \\
Payment & 0.603 & 0.598 & 0.693 & 0.615 & 0.643 & 1.000 \\
\bottomrule
\end{tabular}

\vspace{0.14cm}
\begin{minipage}{0.94\textwidth}
\scriptsize
\justifying
\textbf{用于 embedding 的任务描述示例（入口导航与服务定位）：}
用户刚进入一个数字服务时，需要从首页、菜单、搜索框、分类标签或个人中心中判断应该点击哪里，找到目标功能入口。这个任务的核心困难不是输入信息，而是理解页面结构、识别图标含义、在多个入口之间做选择，并在跳转后确认自己没有走错路径。常见失败包括找不到目标服务、点进错误页面、返回后迷路、页面层级太深。
\end{minipage}

\vfill
{\raggedright\tiny
备注：相似度矩阵用智谱 \texttt{embedding-3} 对任务描述做向量编码，再计算余弦相似度。\par}
\vspace*{-0.12cm}
\end{frame}

% ============================================================
\section{Memory 设计}
% ============================================================

\begin{frame}[t]{复用原生 AgentSociety memory}
\small
\begin{columns}[T]
\begin{column}{0.48\textwidth}
\textbf{原生结构}
\begin{itemize}
    \item \texttt{KVMemory}: key-value 状态
    \item \texttt{StreamMemory}: 时间顺序经历流
    \item \texttt{VectorStore}: 语义检索
    \item \texttt{Memory}: 暴露 \texttt{status} 与 \texttt{stream}
\end{itemize}
\end{column}

\begin{column}{0.48\textwidth}
\textbf{我们复用了}
\begin{itemize}
    \item 写入：\texttt{memory.stream.add}
    \item 检索：\texttt{memory.stream.search}
    \item topic 分类语义
    \item 原生返回文本格式
\end{itemize}
\end{column}
\end{columns}

\vspace{0.18cm}
\scriptsize
\textbf{写入的 description 示例：}\\
\texttt{[digital\_task\_episode][family=payment\_risk\_confirmation]}\\
\texttt{[outcome=failure\_after\_attempt][strategy=try\_self]}\\
\texttt{[help=false][uncontrollability=0.82]}\\
我这次自己尝试支付结算，但还是失败了。页面反复报错，这让我感觉这件事不太受自己控制。\\
\vspace{0.08cm}
\textbf{检索返回时的原生外层格式：}\\
\texttt{- [digital\_task\_episode]: description [day: ..., time: ..., location: ...]}
\end{frame}

\begin{frame}[t]{Stream Memory 接入流程}
\scriptsize
\vspace{1.05cm}
\renewcommand{\arraystretch}{1.18}
\begin{tabular}{@{}L{0.18\linewidth}L{0.30\linewidth}L{0.44\linewidth}@{}}
\toprule
环节 & 做什么 & 接入点 \\
\midrule
写入 & 记录经历流 & 任务、求助、恢复等经历都写进 \texttt{stream}，保持时间顺序和 topic 分类 \\
检索 & 回忆相关经历 & 在 appraisal / attribution / interview 前按 topic 和 query 取回相关内容 \\
审计 & 记录使用痕迹 & retrieval 会带上 condition、status、hash、count 等审计信息 \\
\bottomrule
\end{tabular}
\end{frame}

\end{document}
